\section{Coupling is required for learning}
\label{sec:results}

\subsection{Cross-talk is the difference between a competent detector and a
non-learner}
\label{sec:results-main}
We trained each of the three configurations of Section~\ref{sec:configs} on the
base cued change-detection task (Section~\ref{sec:env-base}) with identical task,
algorithm, hyperparameters and number of updates, and evaluated each on 500 fresh
episodes under a stochastic policy. Outcomes were scored by signal-detection
category --- hit, miss, false alarm, correct rejection --- and summarised by
overall accuracy, hit rate on change trials, false-alarm rate on no-change trials,
and reaction time on hits (Table~\ref{tab:configs}, Figure~\TODO{fig:behavior}).

\begin{table}[t]
\centering
\caption{Behavioural outcomes of the three pathway configurations on the base task
(500 episodes each). The shared-memory (cross-talk) agent is a calibrated detector;
both the independent and swapped-routing agents collapse to never declaring a
change. Numbers are the trained checkpoints reported in the cross-talk analysis;
final-figure values to be regenerated.}
\label{tab:configs}
\small
\begin{tabular}{lcccccl}
\toprule
Configuration & Accuracy & Hit rate & FA rate & RT (frames) & Phenotype \\
\midrule
Shared memory (cross-talk) & \textbf{0.868} & \textbf{0.794} & 0.067 & 1.95 & active, calibrated detector \\
Independent (no cross-talk) & 0.504 & 0.000 & 0.000 & --- & always-wait collapse \\
Swapped routing & 0.420 & 0.000 & 0.000 & --- & always-wait collapse \\
\bottomrule
\end{tabular}
\end{table}

The shared-memory agent is a genuine psychophysical detector: it catches roughly
four of every five real changes, rarely fires on no-change trials
($\mathrm{FAR}=6.7\%$), and responds quickly --- about two frames after change
onset. The independent agent does not partially learn the task; across 500 episodes
it \emph{never once} declares a change. Its hit rate is exactly zero and its
accuracy, $50.4\%$, is precisely the rate obtained by waiting out every trial and
being correct only on the half that contain no change. The policy has retreated to
the single safe action --- always wait --- that can never incur a false-alarm
penalty, and it never discovers that declaring a change is ever worthwhile. The
$50.4$-percentage-point accuracy gap, and the $79.4$-point hit-rate gap, between the
coupled and uncoupled designs is not a matter of degree: it is the difference
between a competent detector and a non-learner, produced by a single structural edit.

\subsection{The collapse is specific to coupling, not to ``having two streams''}
\label{sec:results-swap}
The swapped-routing control rules out the trivial explanation that any split of
policy and value into separate streams fails (or that any such split succeeds).
Configuration~3 \emph{keeps} the cross-talk and the shared memory and changes only
which pathway issues the press decision; it collapses identically --- zero hits,
always wait. Two conditions are therefore jointly necessary. The pathways must be
\emph{coupled} through a shared memory (Configuration~2 fails without it), and the
\emph{grounded, image-reading pathway must drive the policy} (Configuration~3 fails
when it does not). Only the configuration that satisfies both --- the shared-memory
agent with the $\mu$ pathway commanding the actor --- learns. We return to why in
Section~\ref{sec:mechanism}.

\subsection{The working agent reproduces the attentional cueing signature}
\label{sec:results-signature}
A non-collapsed accuracy number is necessary but not sufficient evidence that the
right computation has emerged. The discriminating test is whether the agent
displays the behavioural fingerprint of covert spatial attention: a benefit for
changes at the cued location that \emph{scales with how reliable the cue is}
\citep{posner1980attention, carrasco2011visual, mayo2016graded}. Detailed
characterisation of the shared-memory agent confirms this signature
(Figure~\TODO{fig:cueing}).

\paragraph{Reliability-scaled validity benefit.}
Splitting performance by cue reliability reveals a graded cueing effect: the
advantage for a change at the cued location over an equivalent change at an uncued
location \emph{grows} as the cue's reliability rises from $0.25$ to $1.00$
(\TODO{valid$-$invalid $\Delta$ by reliability}). At the lowest reliability, where
the cue carries no predictive spatial information, there is no cueing effect; at the
highest, a robust advantage emerges. The benefit is expressed primarily as a
leftward shift of the psychometric function --- a lower perceptual threshold for the
attended location --- mirroring the canonical sensitivity-enhancing signature of
attention in primates \citep{treue1999reshaping, reynolds2000attention,
cameron2002sens}.

\paragraph{Value-directed responding.}
Because the cue colour sets the reward magnitude, the agent's behaviour can also be
read for value direction. Across the cue-colour reward structure the agent grades
its responding by value (\TODO{value effect on criterion / response rate}), without
abandoning low-value trials --- the split between the policy and value pathways lets
the value information shape behaviour without collapsing detection. The interaction
between reliability and value, and the relative weight the agent places on each, is
quantified in Section~\TODO{sec:value-analysis}.

\paragraph{Change-locked, fast decisions.}
The agent's declarations are tightly locked to change onset (median reaction time
$\approx 2$ frames; Table~\ref{tab:configs}), and attention allocated to the cued
location is maintained through the blank interval after the cue and reactivated in
anticipation of the change --- the recurrent self-attention dynamics that a
single-stream version of this model was previously shown to reproduce
\citep{herman2017color, ghose2002attentional, nobre2018anticipated}. The two
collapsed variants, by contrast, have no behavioural signature to characterise: an
agent that never presses produces a flat zero hit rate at every reliability and
validity condition. The capacity to express spatial attention is downstream of the
capacity to act at all, and only the coupled architecture clears that first bar.
